% =====================================================================
%  educcion.tex
%  Transcripción del registro de educción del "Reporte de educción e
%  ilación de requisitos de RedPepper", secciones 3.2, 4.2, 5.2 y 6.2.
%  Cada entrada conserva su identificador, contenido, fuente y
%  naturaleza tal como aparecen en el reporte de origen.
% =====================================================================

\section{Educción}
\label{sec:educcion}

\noindent{\small\textbf{Versión 3} \textperiodcentered{} documento de trabajo sujeto a validación con el interesado.}\par\vspace{0.5em}

La educción es la obtención y documentación sistemática de información de
los interesados y del entorno del problema. Su registro conserva qué se
manifestó o describió, quién o qué lo originó, mediante qué técnica se
obtuvo y dónde está su evidencia. Su salida inmediata es un registro de
información del evento, no una decisión técnica ni un requisito aprobado.

Las cuatro tablas siguientes reproducen los registros documentales de los
eventos \textbf{EV-PRE-01} (preentrevista), \textbf{EV-OBS-01}
(observación de campo), \textbf{EV-ENT-01} (entrevista operativa) y
\textbf{EV-ENT-02} (entrevista de profundización). Las entradas son
síntesis del material de origen, tal como el propio reporte advierte; no
constituyen una transcripción literal nueva ni una certificación de una
sesión adicional. El total es de 52 entradas de educción.

\subsection{Evento EV-PRE-01: preentrevista sobre la idea inicial}
\label{subsec:educcion-pre}

La Tabla~\ref{tab:e-pre} reúne las catorce entradas del evento
EV-PRE-01, obtenidas de F-06 (doce solicitudes escritas, aclaración de
emisión SUNAT y modelo mensual).

\begin{longtable}{L{2.05cm} L{6.25cm} L{3.1cm} L{2.85cm}}
\caption{Registro de educción del evento EV-PRE-01 (14 entradas).}
\label{tab:e-pre}\\
\toprule
\textbf{ID} & \textbf{Descripción} & \textbf{Fuente} & \textbf{Observaciones} \\
\midrule
\endfirsthead
\toprule
\textbf{ID} & \textbf{Descripción} & \textbf{Fuente} & \textbf{Observaciones} \\
\midrule
\endhead
\bottomrule
\endfoot
E-PRE-01 & 1. Ventas del día. & F-06, punto 1 & Necesidad declarada. \\
E-PRE-02 & 2. Ingresar los costos directos del producto que se vende. & F-06, punto 2 & Necesidad declarada. \\
E-PRE-03 & 3. Que al final del día me dé el total de ventas del día. & F-06, punto 3 & Necesidad declarada. \\
E-PRE-04 & 4. Que me dé el margen bruto del día que sale de la resta de las ventas totales menos el total del costo directo. & F-06, punto 4 & Relación de cálculo declarada. \\
E-PRE-05 & 5. Las compras se hacen para la semana; el software me debe prorratear las ventas entre lunes y viernes para que el margen bruto diario sea real. Cualquier compra adicional debe sumarse en el día de la compra. & F-06, punto 5 & Regla expresada con ambigüedad. \\
E-PRE-06 & 6. Registrar los costos indirectos del mes: salarios, luz, agua, internet y contabilidad. & F-06, punto 6 & Necesidad declarada. \\
E-PRE-07 & 7. Dar el margen antes de impuestos, BAI, beneficio antes de impuestos e intereses bancarios si los hay. & F-06, punto 7 & Subtotal declarado ambiguo. \\
E-PRE-08 & 8. Calcular los impuestos del mes de acuerdo a las ventas. & F-06, punto 8 & Necesidad declarada incompleta. \\
E-PRE-09 & 9. Dar el margen neto o ganancia del mes. & F-06, punto 9 & Necesidad declarada. \\
E-PRE-10 & 10. Brindar un Estado de Resultados mensual de acuerdo con la nueva NIIF 18. & F-06, punto 10 & Necesidad de reporte y referencia externa. \\
E-PRE-11 & 11. Almacenar facturas de las compras o boletas para justificar ante SUNAT. & F-06, punto 11 & Necesidad de sustento documental. \\
E-PRE-12 & 12. Tener una pestaña de carta con productos y precios para resumir cuáles se venden más y menos y decidir cambios. & F-06, punto 12 & Necesidad funcional y finalidad. \\
E-PRE-13 & Los comprobantes de pago del negocio se sacan de SUNAT. Se solicita un reporte diario de gastos que respalde el informe mensual con los datos pedidos. & F-06, aclaración posterior & Aclaración de alcance. \\
E-PRE-14 & Modelo mensual: ventas menos costos directos igual margen bruto; menos indirectos; margen BAI; menos impuestos e intereses; igual margen neto. El reporte de gastos demuestra los costos directos e indirectos. & F-06, modelo A & Plantilla y relación de sustento. \\
\end{longtable}

\subsection{Evento EV-OBS-01: observación de campo}
\label{subsec:educcion-obs}

La Tabla~\ref{tab:e-obs} reúne las cinco entradas del evento EV-OBS-01,
obtenidas de F-03, sección 2.2, con contraste en F-01, secciones 2 y 5.

\begin{longtable}{L{2.05cm} L{6.25cm} L{3.1cm} L{2.85cm}}
\caption{Registro de educción del evento EV-OBS-01 (5 entradas).}
\label{tab:e-obs}\\
\toprule
\textbf{ID} & \textbf{Descripción} & \textbf{Fuente} & \textbf{Observaciones} \\
\midrule
\endfirsthead
\toprule
\textbf{ID} & \textbf{Descripción} & \textbf{Fuente} & \textbf{Observaciones} \\
\midrule
\endhead
\bottomrule
\endfoot
E-OBS-01 & Se describe una inspección del mostrador y la cocina durante la atención comercial. & F-03, 2.2 & Observación reportada. \\
E-OBS-02 & La observación se reporta como preparación de pedidos al instante, sin productos precocidos. & F-03, 2.2; F-01, 5 & Observación reportada. \\
E-OBS-03 & La preparación genera una espera propia del proceso culinario. & F-03, 2.2 & Observación reportada. \\
E-OBS-04 & El documento describe mostrador, equipos de cocina y área de despacho. & F-03, figura 2 y 2.2 & Contexto físico documentado. \\
E-OBS-05 & El informe propone una interfaz táctil de despacho ágil como resultado de la observación. & F-03, 2.2 & Propuesta técnica documentada. \\
\end{longtable}

\subsection{Evento EV-ENT-01: primera entrevista documentada}
\label{subsec:educcion-ent1}

La Tabla~\ref{tab:e-ent1} reúne las cinco entradas del evento EV-ENT-01,
obtenidas de F-03, sección 2.1 (guion operativo de cuatro preguntas), con
contraste en la carta F-07.

\begin{longtable}{L{2.05cm} L{6.25cm} L{3.1cm} L{2.85cm}}
\caption{Registro de educción del evento EV-ENT-01 (5 entradas).}
\label{tab:e-ent1}\\
\toprule
\textbf{ID} & \textbf{Descripción} & \textbf{Fuente} & \textbf{Observaciones} \\
\midrule
\endfirsthead
\toprule
\textbf{ID} & \textbf{Descripción} & \textbf{Fuente} & \textbf{Observaciones} \\
\midrule
\endhead
\bottomrule
\endfoot
E-ENT1-01 & Jerry explica que redujeron la variedad de platos por costos y rentabilidad. & F-03, 2.1, pregunta 1 & Declaración y motivación de negocio. \\
E-ENT1-02 & La elaboración manual de carnes de hamburguesa elevaba el precio final y los clientes no aceptaban ese precio. & F-03, 2.1, pregunta 1 & Antecedente declarado. \\
E-ENT1-03 & Jerry afirma que no adelantan papas fritas ni hamburguesas y que preparan al recibir el pedido. & F-03, 2.1, pregunta 2 & Regla operativa declarada. \\
E-ENT1-04 & Las papas llegan picadas y el pollo deshilachado se compra a una distribuidora denominada Don Pollo. & F-03, 2.1, pregunta 3 & Dato de abastecimiento declarado. \\
E-ENT1-05 & Las gaseosas se compran a distribuidores; los refrescos naturales usan pulpa y agua hervida o mineral. & F-03, 2.1, pregunta 4 & Descripción declarada de producto y preparación. \\
\end{longtable}

\subsection{Evento EV-ENT-02: entrevista de profundización}
\label{subsec:educcion-ent2}

La Tabla~\ref{tab:e-ent2} reúne las veintiocho entradas del evento
EV-ENT-02, correspondientes a las preguntas P-01 a P-28 de F-08
(Preguntas.pdf). Las respuestas se resumen preservando objeciones,
matices y temas no resueltos.

{\small
\begin{longtable}{L{2.05cm} L{6.25cm} L{3.1cm} L{2.85cm}}
\caption{Registro de educción del evento EV-ENT-02 (28 entradas).}
\label{tab:e-ent2}\\
\toprule
\textbf{ID} & \textbf{Descripción} & \textbf{Fuente} & \textbf{Observaciones} \\
\midrule
\endfirsthead
\toprule
\textbf{ID} & \textbf{Descripción} & \textbf{Fuente} & \textbf{Observaciones} \\
\midrule
\endhead
\bottomrule
\endfoot
E-ENT2-01 & P-01: el programa debe abrirse en el negocio y fuera; Jerry revisa desde casa. & F-08, P-01 & Necesidad y contexto de uso. \\
E-ENT2-02 & P-02: en el local se dedica a la operación; administra desde su escritorio en casa. & F-08, P-02 & Contexto declarado. \\
E-ENT2-03 & P-03: acceso del dueño, contador, administrador y a veces chef por stocks. & F-08, P-03 & Necesidad de usuarios. \\
E-ENT2-04 & P-04: pantalla táctil para servicio y mouse para administración. & F-08, P-04 & Preferencia de interacción. \\
E-ENT2-05 & P-05: al indicar enter quiere obtener todas las ventas del día. & F-08, P-05 & Necesidad de resumen. \\
E-ENT2-06 & P-06: seleccionar salchipapa y cremas en pantalla; sale una comanda para el chef. & F-08, P-06 & Necesidad funcional. \\
E-ENT2-07 & P-07: reportes diarios, semanales y mensuales; el usuario decide inicio y fin de jornada. & F-08, P-07 & Necesidad y regla declarada. \\
E-ENT2-08 & P-08: todo sistema debe guardar historial. & F-08, P-08 & Necesidad declarada. \\
E-ENT2-09 & P-09: puede borrar antes de procesar; si ya vendió el tratamiento es otro. & F-08, P-09 & Distinción declarada incompleta. \\
E-ENT2-10 & P-10: considera importante buscar día o plato en historial. & F-08, P-10 & Necesidad funcional. \\
E-ENT2-11 & P-11: al teclear el producto agotado no sale y eso le alerta. & F-08, P-11 & Comportamiento esperado. \\
E-ENT2-12 & P-12: el usuario dispone si la compra se consume en el día o en la semana. & F-08, P-12 & Regla declarada. \\
E-ENT2-13 & P-13: responde sí al cargo del día, pero añade según montos y potestad del usuario. & F-08, P-13 & Respuesta condicionada. \\
E-ENT2-14 & P-14: repartir entre días hábiles; artículos personales se compran aparte. & F-08, P-14 & Dos reglas declaradas. \\
E-ENT2-15 & P-15: si sube el costo, Jerry actualiza el precio para proteger el margen. & F-08, P-15 & Política comercial declarada. \\
E-ENT2-16 & P-16: acepta corregir montos; supervisa el estado de resultados junto al contador. & F-08, P-16 & Necesidad de corrección. \\
E-ENT2-17 & P-17: acepta plantilla; ventas menos directos, margen menos indirectos y luego impuestos. & F-08, P-17 & Necesidad y relación de cálculo. \\
E-ENT2-18 & P-18: solicita NIIF 18; contador propone y Jerry autoriza. & F-08, P-18 & Intención y regla de aprobación. \\
E-ENT2-19 & P-19: espera que el cambio se aplique cuando SUNAT indique la fecha. & F-08, P-19 & Expectativa sobre vigencia. \\
E-ENT2-20 & P-20: acepta estructura formal aprobada por el contador. & F-08, P-20 & Necesidad de validación. \\
E-ENT2-21 & P-21: no quiere explicación paso a paso; desea estado de resultados mensual para analizar. & F-08, P-21 & Respuesta negativa y preferencia. \\
E-ENT2-22 & P-22: el contador tiene facturas electrónicas; desde el gasto desea hallar factura, fecha y concepto. & F-08, P-22 & Necesidad de trazabilidad. \\
E-ENT2-23 & P-23: menciona dualidad y afirma que Camaleón está homologado con SUNAT. & F-08, P-23 & Opinión y referencia a otro producto. \\
E-ENT2-24 & P-24: acepta bloquear anulación directa y exigir nota de ajuste autorizada. & F-08, P-24 & Regla declarada. \\
E-ENT2-25 & P-25: acepta compresión automática de PDF pesado. & F-08, P-25 & Necesidad de almacenamiento. \\
E-ENT2-26 & P-26: acepta avisos por números negativos y fechas incorrectas. & F-08, P-26 & Necesidad de validación. \\
E-ENT2-27 & P-27: acepta que no emite automáticamente, pero opina que SUNAT querrá afiliación. & F-08, P-27 & Acuerdo con reserva. \\
E-ENT2-28 & P-28: insiste en que cualquier pedido que llegue sigue siendo venta y se registra. & F-08, P-28 & Necesidad comercial. \\
\end{longtable}
}

\paragraph{Salidas de la educción.} El reporte de origen identifica cuatro
registros documentales: EV-PRE-01-ED, EV-OBS-01-ED, EV-ENT-01-ED y
EV-ENT-02-ED. Ninguno contiene todavía decisiones de implementación
aprobadas. F-05 (audio de WhatsApp) se conserva como fuente adicional
separada de los eventos y no se utiliza como premisa hasta contar con una
transcripción revisable y localizadores temporales.
